% !TeX root = ../main.tex
% 7. fejezet — Vadásztársaságok és tagságkezelés
% Hivatkozott források: Backend/HUNting.API/Controllers/SocietiesController.cs, NewsController.cs,
%   Services/SocietyService.cs, Models/Societies/*, Data/HuntingContext.cs, Data/SeedData.cs,
%   Frontend/HUNting/src/services/societiesApi.ts, pages/Societies/SocietiesPage.tsx,
%   SocietyDetailPage.tsx, SocietyManagementPage.tsx, MySocietiesPage.tsx

\chapter{Vadásztársaságok és tagságkezelés}
\label{chap:tarsasagok}

A vadásztársasági funkciócsoport a HUNting közösségi magja: itt találkozik a két fióktípus. Az
egyéni vadászok társaságokat kereshetnek és tagfelvételüket kérhetik, a társasági fiókok pedig a
kérelmek elbírálásával, a taglistával és a belső bejegyzésekkel igazgatják közösségüket. A fejezet
a tagfelvételi kérelem életciklusára mint a rendszer legjellegzetesebb üzleti folyamatára
összpontosít, zárásként pedig tisztázza a hírfolyam (NewsFeed) jelenlegi, részlegesen megvalósított
állapotát.

\section{A funkció célja és felhasználói működése}
\label{sec:soc-cel}

A társaságkatalógus nyilvános: a látogató név szerint szűrhető listában böngészheti a regisztrált
társaságokat, és megnyithatja bármelyik profilját, amely a bemutatkozást, az elérhetőségeket, az
alapítás dátumát és a taglétszámot mutatja. Bejelentkezett, egyéni fiókkal rendelkező felhasználó a
profiloldalról tagfelvételi kérelmet nyújthat be, opcionális bemutatkozó üzenettel. A kérelemről a
társaság tulajdonosa dönt a társaságkezelő felületen, ahol a függőben lévő jelentkezések, a
taglista és a belső posztok is egy helyen érhetők el. Elfogadás után az új tag a „Társaságaim”
oldalon látja tagságait, és társaságot választva olvashatja annak belső bejegyzéseit -- a posztok
tehát csak a tagság megszerzése után válnak láthatóvá.

% TODO: SC-05 és SC-06 képernyőkép (figures/sc-05-societies.png, sc-06-society-detail.png)
\begin{figure}[htbp]
	\centering
	\fbox{\parbox[c][5cm][c]{0.85\textwidth}{\centering SC-05/SC-06: A társaságkatalógus és a társasági profil a jelentkezési űrlappal (elkészítendő)}}
	\caption{A társaságkatalógus és egy társaság profilja a tagfelvételi űrlappal}
	\label{fig:societies-screens}
\end{figure}

\section{Adatmodell}
\label{sec:soc-adatmodell}

A funkciócsoport öt entitásra épül (\texttt{Models/Societies/}). A \texttt{Society} a társaság
törzsadatait tárolja (egyedi indexszel védett e-mail-cím, név, telefonszám, bemutatkozás,
létrehozási idő), és a~\ref{chap:felhasznalok}.~fejezetben bemutatott módon egy-egy kapcsolatban
áll tulajdonos felhasználójával. A társaságon belüli rangokat a \texttt{SocietyRole} szótártábla
írja le; a kezdeti adatfeltöltés négy szerepet hoz létre (Vezeto, Gazda, Vador, Tag).

A tagságot a \texttt{SocietyMember} kapcsolótábla valósítja meg, amely a klasszikus több-a-többhöz
mintát követi, de a puszta kapcsolaton túl attribútumokat is hordoz: a tag társaságon belüli
szerepét és a csatlakozás időpontját. Elsődleges kulcsa összetett -- a
\texttt{(SocietyId, UserId)} pár --, ami adatbázisszinten zárja ki, hogy ugyanaz a felhasználó
kétszer legyen tagja ugyanannak a társaságnak (\ref{lst:member-config}.~kódrészlet). Ugyanilyen
összetett kulcsú a \texttt{SocietyMembershipRequest} entitás is, vagyis egy felhasználóhoz
társaságonként legfeljebb egy kérelemrekord tartozhat; az ebből következő „újranyitási” viselkedést
a~\ref{sec:soc-backend}.~alfejezet mutatja be. A kérelem állapotát a
\texttt{MembershipRequestStatus} felsorolt típus írja le (\texttt{Pending}, \texttt{Approved},
\texttt{Rejected}), emellett a rekord az elbírálás időpontját és az elbíráló azonosítóját is
rögzíti.

\begin{lstlisting}[language={[Sharp]C}, caption={A tagsági kapcsolótábla konfigurációja a \texttt{HuntingContext} osztályban (kivonat)}, label={lst:member-config}]
modelBuilder.Entity<SocietyMember>()
    .HasKey(x => new { x.SocietyId, x.UserId });

modelBuilder.Entity<SocietyMember>()
    .HasOne(x => x.Society)
    .WithMany(x => x.Members)
    .HasForeignKey(x => x.SocietyId)
    .OnDelete(DeleteBehavior.Cascade);

modelBuilder.Entity<SocietyMember>()
    .HasOne(x => x.User)
    .WithMany(x => x.Societies)
    .HasForeignKey(x => x.UserId)
    .OnDelete(DeleteBehavior.Cascade);

modelBuilder.Entity<SocietyMember>()
    .HasOne(x => x.Role)
    .WithMany(x => x.Members)
    .HasForeignKey(x => x.RoleId)
    .OnDelete(DeleteBehavior.Restrict);
\end{lstlisting}

A törlési szabályok a rekord jellegéhez igazodnak: a tagság a társasággal és a felhasználóval
együtt kaszkádolva törlődik, a szerephivatkozás viszont \texttt{Restrict}, így használatban lévő
társasági szerep nem távolítható el. A kérelemnél árnyaltabb a kép: a társaság törlése a
kérelmeket is viszi, a kérelmező felhasználó törlése tiltott, amíg kérelme létezik
(\texttt{Restrict}), az elbíráló törlésekor pedig a hivatkozás \texttt{SetNull} értékre vált --
a döntés ténye tehát az elbíráló fiókjának megszűnése után is megmarad. A belső bejegyzéseket a
\texttt{SocietyPost} entitás tárolja (cím, tartalom, opcionális képútvonal, szerző, időbélyeg),
amely a társasággal és a szerzővel együtt kaszkádolva törlődik.

% TODO: E-02 osztálydiagram (figures/e-02-society-domain.pdf) — ApplicationUser, Society,
%       SocietyMember, SocietyMembershipRequest, SocietyRole, SocietyPost kapcsolatokkal.
\begin{figure}[htbp]
	\centering
	\fbox{\parbox[c][6.5cm][c]{0.9\textwidth}{\centering E-02: A társasági domain UML-osztálydiagramja (elkészítendő)}}
	\caption{A társasági domain entitásai és kapcsolataik}
	\label{fig:society-domain}
\end{figure}

\section{Backend-megvalósítás}
\label{sec:soc-backend}

A funkciót a \texttt{SocietiesController} és a \texttt{SocietyService} párosa valósítja meg. A
szolgáltatás minden adatelérési műveletét közvetlenül a \texttt{HuntingContext}-en végzi: a
társaság-törzsadatok egyszerű CRUD-műveleteit éppúgy, mint az összetettebb lekérdezéseket és a
kapcsolótáblák kezelését. A végpontokat
a~\ref{tab:soc-endpoints}.~táblázat összegzi.

\begin{table}[htbp]
	\centering
	\caption{A \texttt{SocietiesController} végpontjai}
	\label{tab:soc-endpoints}
	\begin{tabular}{|l|l|p{5.2cm}|}
		\hline
		\textbf{Végpont} & \textbf{Hozzáférés} & \textbf{Feladat} \\
		\hline
		GET \texttt{/api/Societies} & anonim & Társaságok listája \\
		\hline
		GET \texttt{/api/Societies/\{id\}} & anonim & Egy társaság részletei \\
		\hline
		GET \texttt{/api/Societies/my} & bejelentkezett & A hívó tagságai \\
		\hline
		GET \texttt{/api/Societies/\{id\}/members} & tulajdonos v. admin & Taglista \\
		\hline
		POST \texttt{/api/Societies/\{id\}/join} & bejelentkezett Person & Tagfelvételi kérelem \\
		\hline
		GET \texttt{/api/Societies/\{id\}/requests} & tulajdonos v. admin & Függő kérelmek \\
		\hline
		POST \texttt{.../requests/\{userId\}/approve} & tulajdonos v. admin & Kérelem elfogadása \\
		\hline
		POST \texttt{.../requests/\{userId\}/reject} & tulajdonos v. admin & Kérelem elutasítása \\
		\hline
		GET \texttt{/api/Societies/\{id\}/posts} & tag, tulajdonos v. admin & Belső posztok \\
		\hline
		POST \texttt{/api/Societies/\{id\}/posts} & tulajdonos v. admin & Poszt létrehozása képpel \\
		\hline
		DELETE \texttt{.../posts/\{postId\}} & tulajdonos v. admin & Poszt törlése \\
		\hline
	\end{tabular}
\end{table}

\subsection{A tagfelvételi kérelem életciklusa}

A jelentkezési végpont először a fióktípust ellenőrzi: a token \texttt{account\_type} claimjéből
megállapítja, hogy a hívó társasági fiók-e, és ha igen, elutasítja a kérést -- társaság nem
csatlakozhat másik társasághoz. Ezután a szolgáltatás a kérelem állapotgépét hajtja végre
(\ref{lst:join-lifecycle}.~kódrészlet): ha a hívó már tag, a kérés hibával zárul; ha van függőben
lévő kérelme, szintén; ha viszont van korábbi -- elfogadott vagy elutasított -- kérelemrekordja,
azt a rendszer \emph{újranyitja}: az állapotot \texttt{Pending} értékre állítja, frissíti az
üzenetet és a beadás időpontját, az elbírálási mezőket pedig kiüríti. Új rekord csak akkor jön
létre, ha a felhasználónak még egyáltalán nem volt kérelme az adott társaságnál -- ezt a
megoldást az összetett kulcs kényszeríti ki, amely társaságonként és felhasználónként egyetlen
kérelemrekordot enged.

\begin{lstlisting}[language={[Sharp]C}, caption={A kérelem létrehozása és újranyitása a \texttt{SocietyService} osztályban (kivonat)}, label={lst:join-lifecycle}]
var existingMember = await _ctx.SocietyMembers
    .AnyAsync(m => m.SocietyId == societyId && m.UserId == userId);
if (existingMember)
{
    throw new InvalidOperationException("Mar tagja vagy a tarsasagnak.");
}

var request = await _ctx.SocietyMembershipRequests.FindAsync(societyId, userId);
if (request != null)
{
    if (request.Status == MembershipRequestStatus.Pending)
    {
        throw new InvalidOperationException("Mar van fuggoben levo kerelmed.");
    }

    request.Status = MembershipRequestStatus.Pending;
    request.Message = message;
    request.CreatedAt = DateTime.UtcNow;
    request.ProcessedAt = null;
    request.ProcessedByUserId = null;
}
else
{
    request = new SocietyMembershipRequest { /* ... */ };
    _ctx.SocietyMembershipRequests.Add(request);
}
\end{lstlisting}

Az elbírálás a tulajdonos vagy adminisztrátor kiváltsága, amit a kontroller közös
\texttt{IsOwnerOrAdmin} segédmetódusa ellenőriz. Az elfogadó művelet idempotens: már elfogadott
kérelemre nem fut le újra. Sikeres elfogadáskor a kérelem állapota \texttt{Approved} lesz, a
rekord rögzíti az elbírálót és az időpontot, majd -- ha a felhasználó időközben nem lett tag --
létrejön a tagsági rekord az alapértelmezett „Tag” társasági szereppel (ennek hiányában az első
elérhető szereppel). Az elutasítás csak az állapotot és az elbírálási mezőket írja; a korábban
bemutatott újranyitási szabály miatt az elutasított jelentkező később újra próbálkozhat. A teljes
életciklust a~\ref{fig:request-lifecycle}.~ábra foglalja össze.

% TODO: F-02 állapotdiagram (figures/f-02-request-lifecycle.pdf) — (nincs) → Pending →
%       Approved / Rejected → újranyitás Pendingre.
\begin{figure}[htbp]
	\centering
	\fbox{\parbox[c][5.5cm][c]{0.9\textwidth}{\centering F-02: A tagfelvételi kérelem állapotdiagramja (elkészítendő)}}
	\caption{A tagfelvételi kérelem életciklusa}
	\label{fig:request-lifecycle}
\end{figure}

\subsection{Belső bejegyzések}

A posztok olvasási joga szélesebb, mint az írási: a listát a társaság bármely tagja, a tulajdonos
és az adminisztrátor kérheti le -- a kontroller a tagságot és a tulajdonlást közvetlen
adatbázis-lekérdezéssel ellenőrzi --, új posztot viszont csak a tulajdonos vagy adminisztrátor
hozhat létre. A poszt többrészes kérésben érkezik (cím, tartalom, opcionális kép, 10~MiB
kéréslimittel), a képkezelés pedig a hirdetéseknél megismert mintát követi:
kiterjesztés-fehérlista, GUID fájlnév, \texttt{uploads/societies} célkönyvtár. Poszt törlésekor a
kontroller ellenőrzi, hogy a bejegyzés valóban az útvonalban megadott társasághoz tartozik-e,
majd a képfájlt és a rekordot is eltávolítja.

\section{Frontend-megvalósítás}
\label{sec:soc-frontend}

\subsection{Katalógus és jelentkezés}

A katalógusoldal (\texttt{SocietiesPage}) a teljes társaságlistát kéri le, és a névszűrést --
a lista mérete miatt elfogadható kompromisszumként -- kliensoldalon végzi. A profiloldal
(\texttt{SocietyDetailPage}) a \texttt{["society", societyId]} kulcson kérdezi le az adatokat, és
a jelentkezési űrlapot csak akkor jeleníti meg, ha a látogató bejelentkezett és fióktípusa
\texttt{Person} -- ugyanaz a szabály, amelyet a backend a token claimjéből kikényszerít. A
jelentkezés mutációja sikeres beküldés után visszaigazoló üzenetet jelenít meg és üríti az
üzenetmezőt; a szerver által dobott üzleti hibák (már tag, már van függő kérelem) az űrlap
hibaszövegeként jelennek meg.

\subsection{Társaságkezelő és tagsági nézet}

A társaságkezelő oldal (\texttt{SocietyManagementPage}) hozzáférési feltétele kliensoldalon is
kettős: bejelentkezett, \texttt{Society} típusú fiók kell, amelyhez tartozik birtokolt társaság
(\texttt{user.ownedSociety}) -- ennek hiányában az oldal magyarázó üzenetet jelenít meg. A
társaság azonosítóját a tárolt profil \texttt{ownedSociety} mezőjéből veszi, és három párhuzamos
lekérdezést indít: tagok (\texttt{["society-members", id]}), függő kérelmek
(\texttt{["society-requests", id]}) és posztok (\texttt{["society-posts", id]}). A mutációk
célzottan érvénytelenítik az érintett gyorsítótárakat: az elfogadás a kérelmek \emph{és} a tagok
listáját is frissíti -- hiszen mindkettőt megváltoztatja --, az elutasítás csak a kérelmekét, a
posztműveletek pedig a posztokét. Az új poszt űrlapja cím- és tartalommezőből, valamint natív
fájlválasztóból áll; beküldés után a mezők kiürülnek.

A tagsági nézet (\texttt{MySocietiesPage}) az egyéni fiókok tükörképe: a \texttt{["my-societies"]}
lekérdezés a hívó tagságait adja vissza, a kiválasztott társaság posztjait pedig egy feltételesen
engedélyezett lekérdezés (\texttt{enabled: Boolean(selectedSociety)}) tölti be -- így hálózati
kérés csak akkor indul, amikor a felhasználó ténylegesen társaságot választott. A posztok képei a
\texttt{resolveAssetUrl} segédfüggvényen keresztül kapják meg a backend abszolút címét.

% TODO: SC-07 és SC-08 képernyőkép (figures/sc-07-society-manage.png, sc-08-my-societies.png)
\begin{figure}[htbp]
	\centering
	\fbox{\parbox[c][5cm][c]{0.85\textwidth}{\centering SC-07/SC-08: A társaságkezelő és a „Társaságaim” oldal (elkészítendő)}}
	\caption{A társaságkezelő felület a függő kérelmekkel, valamint a tagsági nézet a belső posztokkal}
	\label{fig:society-manage-screens}
\end{figure}

\section{A hírfolyam (NewsFeed) állapota és a funkciócsoport korlátai}
\label{sec:soc-news}

A kódbázis a belső posztok mellett egy második, hasonló rendeltetésű entitáscsaládot is tartalmaz:
a \texttt{NewsFeed} hírbejegyzést és a hozzá tartozó \texttt{NewsVisibility} láthatósági
szótártáblát (\texttt{MEMBERS}, \texttt{PUBLIC}, \texttt{OFFICERS} kódokkal), két mintahírrel a
kezdeti adatfeltöltésben. Fontos azonban őszintén rögzíteni a megvalósultság tényleges szintjét:
a \texttt{NewsController} teljes CRUD-ot kínál -- az író végpontjait (létrehozás, módosítás,
törlés) az \texttt{[Authorize]} attribútum adminisztrátori szerep-megkötése védi, az olvasó
végpontok anonim hívhatók --, a tárolt láthatósági kód szerinti
szűrés azonban nem valósul meg, a frontenden pedig a hírekhez semmilyen oldal vagy API-hívás nem
tartozik. A hírfolyam tehát \emph{backend-oldali előkészítésnek}
tekintendő, nem működő felhasználói funkciónak; a dolgozat a tagok felé irányuló tényleges
kommunikációs csatornaként a \texttt{SocietyPost} bejegyzéseket dokumentálja. A hiányzó
láthatóságszűrést és felhasználói felületet a továbbfejlesztési irányok között
(\ref{chap:osszefoglalas}.~fejezet) kell rendezni.

A funkciócsoport további ismert korlátai: a katalógusoldal kiemelt „Jelentkezés tagságra” gombjához
jelenleg nem tartozik eseménykezelő -- a működő jelentkezési űrlap a profiloldalon él --; a
taglistáról nincs tageltávolító felületi művelet (a szolgáltatásban a metódus elkészült, végpont
nem tartozik hozzá); a társasági szerepek (Gazda, Vador) pedig a tagrekordban tárolódnak ugyan, de
jogosultsági következményük jelenleg nincs -- a döntéseket minden ponton a tulajdonos-vagy-admin
szabály vezérli.
